\documentclass[10pt,twocolumn]{article}
\usepackage[a4paper,margin=16mm,columnsep=7mm]{geometry}
\usepackage[T1]{fontenc}
\usepackage{lmodern,amsmath,amssymb,amsthm,mathtools,microtype,needspace}
\usepackage{hyperref}
\hypersetup{colorlinks=true,linkcolor=black,citecolor=black,urlcolor=blue}
\newtheorem{theorem}{Theorem}
\newtheorem{proposition}{Proposition}
\newtheorem{lemma}{Lemma}
\title{Kernel-Concavity Obstructions and Fourier-Side Laguerre Bounds\\for the Riemann $\Xi$ Programme}
\author{Tony Newton\\Newton Astro Labs\\United Kingdom\\\texttt{tony.newton79@gmail.com}}
\date{}
\begin{document}
\maketitle
\noindent\textit{Reserved DOI, not yet published:}
\href{https://doi.org/10.5281/zenodo.23461586}{10.5281/zenodo.23461586}.
\begin{abstract}
We give an exact integer-interval certificate showing that the fifth
nonlinear differential concavity iterate of the genuine Riemann
theta kernel is negative at the modular point, while iterates zero
through four are positive. We exhibit two strongly log-concave
Gaussian-polynomial kernels with positive iterates of every order
but Fourier transforms with nonreal zeros. The smaller example
violates the first Fourier-side Laguerre inequality at an explicit
real point. Lean certifies finite algebraic and moment-positivity subclaims, while
an independently audited rational Python interval proves the local
fifth-iterate bound for the actual theta series; full end-to-end
Lean formalization remains open. These results do not establish
or refute the Riemann hypothesis. Priority remains unverified.
\end{abstract}
\section{A fifth-iterate obstruction for the theta kernel}
Set $\mathcal Tf=(f')^2-ff''$ and
\begin{align*}
\Phi(r)&=\sum_{m\ge1}q_m(2q_me^{4r}-3)e^{5r-q_me^{4r}},\\
q_m&=\pi m^2,\qquad s(t)=\Phi(\sqrt t).
\end{align*}
The theta-functional equation makes $\Phi$ even; thus $s$ is analytic
near zero. This is the kernel normalization used by
Planat--Sol\'e~\cite{PS}; constant and argument rescalings do not
change the signs under consideration.

\begin{theorem}[Certified obstruction]\label{thm:five}
One has
\[
-4.813\cdot 10^{65}<\mathcal T^5s(0)
<-4.812\cdot10^{65}
\]
and $\mathcal T^js(0)>0$ for $0\le j\le4$. In particular
$\mathcal T^4s$ is positive but strictly log-convex at zero,
and $\mathcal T^5s<0$ on some punctured interval to the right.
\end{theorem}
The second-level concavity theorem of Planat--Sol\'e~\cite{PS}
is not contradicted by this higher-order local obstruction.

\begin{lemma}[Finite theta tail]
For each $0\le k\le10$, truncation at $m=12$ in the exact
Taylor-coefficient formula
\begin{align*}
a_k&=[t^k]s(t)\\
&=\frac1{(2k)!}\sum_{m\ge1}q_me^{-q_m}P_{2k}(q_m),\\
P_0(x)&=2x-3,\\
P_{j+1}(x)&=(5-4x)P_j(x)+4xP_j'(x),
\end{align*}
has absolute error below $4\cdot10^{-83}$.
\end{lemma}
\begin{proof}
Put $C_j=\sum|\operatorname{coeff}(P_j)|$.
For $x\ge1$, $|P_{2k}(x)|\le C_{2k}x^{2k+1}$.
Using $3<\pi<22/7$ and $e>2$, the omitted $m$-th
summand has magnitude at most
\[
\frac{C_{2k}}{(2k)!}
\left(\frac{22m^2}7\right)^{2k+2}2^{-3m^2}.
\]
The successive ratio for $m\ge13$ and $0\le k\le10$ is
bounded by $(14/13)^{44}2^{-81}<1/2$.
Thus the absolute tail is at most
\[
B_k=\frac{2C_{2k}}{(2k)!}
\left(\frac{22\cdot169}7\right)^{2k+2}2^{-507}.
\]
Exact integer comparisons in the supplied Python certificate verify
$B_k<4\cdot10^{-83}$ for all eleven values of $k$.
\end{proof}

\begin{proof}[Proof of Theorem~\ref{thm:five}]
Only $a_0,\ldots,a_{10}$ are required. If $f=\sum b_kt^k$ then
\begin{align*}
[t^k]\mathcal Tf
&=\sum_{i+j=k}(i+1)(j+1)b_{i+1}b_{j+1}\\
&\quad-\sum_{i+j=k}(j+1)(j+2)b_ib_{j+2}.
\end{align*}
Five successive applications consume exactly ten trailing Taylor
coefficients. To certify these initial coefficients, enclose
$\pi=16\arctan(1/5)-4\arctan(1/239)$ using alternating rational
series with 100 and 40 terms. Enclose $e^{-q_m}$ for $1\le m\le12$
via the reciprocal of $[e^{q_m/1024}]^{1024}$: Taylor degree 80
for the reduced argument has a rigorous omitted bound
$2\cdot2^{-81}/81!<10^{-110}$.
Every arithmetic operation is an outward-rounded integer
operation on intervals with denominator $10^{110}$. Expand
the input coefficient intervals by the analytic theta-tail bound
and propagate them five times by the displayed recurrence.
The source \texttt{certificates/xi\_fifth\_iterate\_exact.py}
reproduces the stated enclosure without third-party packages
or floating-point decisions. The signs for levels $0$ to $4$ are
strictly positive, and level $5$ strictly negative.
Finally $(\log\mathcal T^4s)''(0)
=-\mathcal T^5s(0)/(\mathcal T^4s(0))^2>0$, and analyticity
gives the asserted neighborhood.
\end{proof}

\Needspace{10\baselineskip}
\section{Two all-iteration counterexamples}

\begin{proposition}[Small two-cycle]\label{prop:small}
Let $K(u)=e^{-u^2/4}(u^4+22u^2+241)$ and
$g(t)=K(\sqrt t)$. Put $A(t)=t^2+22t+241$
and $B(t)=t^2+22t+1$. Then, for all $n\ge0$
and $t\ge0$,
\[
\mathcal T^ng(t)=2^{2^n-1}e^{-2^nt/4}
\begin{cases}A(t),&n\text{ even},\\B(t),&n\text{ odd}\end{cases}>0.
\]
Moreover $(\log K)''<-1/4$ for every real $u$, yet
\[
\frac{1}{2\sqrt\pi}\int_{\mathbb R}K(u)e^{izu}du
=e^{-z^2}\big((4z^2-17)^2+8\big)
\]
has four nonreal zeros.
\end{proposition}
\begin{proof}
Direct differentiation shows
$\mathcal TA=2B$ and $\mathcal TB=2A$. Since
$\mathcal T(ce^{-at}P)=c^2e^{-2at}\mathcal TP$,
the iterate formula follows inductively.
With $Q(u)=u^4+22u^2+241$, one has
\begin{align*}
\frac{Q^2}4-QQ''+(Q')^2
&=\frac{u^8}4+15u^6+\frac{571u^4}2\\
&\quad+727u^2+\frac{15665}4>0.
\end{align*}
Thus $(\log K)''=-1/4-
\bigl(Q^2/4-QQ''+(Q')^2\bigr)/Q^2<-1/4$.
Gaussian Fourier differentiation gives the claimed transform;
its zeros satisfy $z^2=(17\pm2\sqrt2 i)/4$.
\end{proof}

\begin{proposition}[Original rational-root construction]
Let
\[
K_*(u)=e^{-u^2/4}(62161+1848u^2+16u^4).
\]
Set $A_*(t)=62161+1848t+16t^2$ and
$B_*(t)=44561+1848t+16t^2$.
For $g_*(t)=K_*(\sqrt t)$,
\[
\mathcal T^n g_*(t)=32^{2^n-1}e^{-2^nt/4}
\begin{cases}A_*(t),&n\text{ even},\\
B_*(t),&n\text{ odd}\end{cases}>0.
\]
The kernel satisfies $(\log K_*)''<-3/8$ on $\mathbb R$,
but its normalized Fourier transform is
\[
F_*(z)=e^{-z^2}\big[(16z^2-255)^2+1024\big],
\]
whose zeros are exactly $z=\pm4\pm i/4$.
\end{proposition}
\begin{proof}
One checks $\mathcal TA_*=32B_*$ and
$\mathcal TB_*=32A_*$. The exponential scaling identity
used above proves the iteration formula. With
$Q_*(u)=62161+1848u^2+16u^4$,
\[
Q_*(u)/8-Q_*''(u)=32593/8+39u^2+2u^4>0
\]
implies the stated log-concavity bound.
Fourier differentiation gives $F_*$; since
$16(4+i/4)^2-255=32i$, it has the claimed root,
and evenness and complex conjugation supply the other three.
\end{proof}

The idea of counterexamples based on positive
Gaussian-polynomial kernels predates this note; see
Csordas~\cite{Cs}, Example 3.12. No priority claim is made
for that general construction.

\section{Fourier-side generalized Laguerre diagnostics}
For the smaller counterexample in Proposition~\ref{prop:small},
write $F(x)=e^{-x^2}H(x)$ with
$H(x)=(4x^2-17)^2+8$. At $x=2$,
\[
H(2)=9,\qquad H'(2)=-32,\qquad H''(2)=496.
\]
Consequently the \emph{transform-side} first Laguerre expression
is strictly negative:
\begin{equation}\label{eq:lag}
(F'(2))^2-F(2)F''(2)
=e^{-8}(1024-4464+162)
=-3278e^{-8}<0.
\end{equation}
This is compatible with positivity of \emph{every} iterated
concavity expression on $g(t)$, and proves that kernel and
transform inequalities are different hierarchies.

For an even real Schwartz kernel and
$F(x)=c\int_{\mathbb R}K(u)e^{ixu}du$, $c>0$, one has
\begin{align}\label{eq:correlation}
(F'(x))^2-F(x)F''(x)
=\frac{c^2}{2}\iint_{\mathbb R^2}
&(u-v)^2K(u)K(v)\\[-0.8em]
&\cdot e^{ix(u+v)}\,du\,dv.\nonumber
\end{align}
Writing $w=u+v$, this is the Fourier transform of the
pointwise nonnegative correlation
\[
G(w)=\frac{c^2}{2}\int_{\mathbb R}
(2u-w)^2K(u)K(w-u)\,du.
\]
However $G(w)\ge0$ does \emph{not} imply its Fourier transform
is nonnegative, as~\eqref{eq:lag} explicitly demonstrates.

For a real entire function whose Taylor expansions permit the
interchange, define the higher generalized Laguerre expressions
through
\[
F(x+iy)F(x-iy)=\sum_{n\ge0}L_n[F](x)y^{2n},
\]
where
\[
L_n[F](x)=\frac{(-1)^n}{(2n)!}
\sum_{k=0}^{2n}(-1)^k{2n\choose k}
F^{(k)}(x)F^{(2n-k)}(x).
\]
Proper transform-level certificates are a more plausible approach
to controlling complex zeros than indefinite kernel iteration.
Wronskian criteria and related approaches are discussed by
Dimitrov--Xu~\cite{DX}.

\section{Machine verification and limitations}
The repository contains \texttt{XiKernel/Algebra.lean}, with
polynomial two-cycle identities and positivity sublemmas,
a complex polynomial-zero identity and exact Laguerre-sign
arithmetic. The Lean files do \emph{not} formalize the analytic
theta-tail lemma, the interval implementation, or RH.
The accompanying Python program implements exact fixed-point
interval arithmetic; its analytic majorants are displayed in
the proof above. A separate polynomial test gives an
independent rational-arithmetic replay of the Gaussian identities.

The repository has a GitHub Actions workflow for Lean compilation.
The GitHub Actions Lean job completed successfully on 19 October
2026 with Lean 4.24.0 and Mathlib (commit \texttt{922ed6c6}).
This certifies the declarations in the algebraic module only:
the analytic theta-tail proof and interval certificate are
independently checked in Python, not Lean.

\section{Interpretation}
Kernel differential iteration, infinite log-concavity of the
\emph{Taylor coefficient sequence}, and generalized Laguerre
inequalities on the \emph{Fourier transform} are distinct.
In particular, Liu--Mao~\cite{LM} prove coefficient
log-concavity without proving RH, and Planat--Sol\'e~\cite{PS}
prove the second-level theta-kernel theorem without asserting RH.

The fifth-iterate sign rules out an infinite differential
concavity strengthening for the \emph{actual theta kernel}.
It neither proves a nonreal zero of the Riemann $\Xi$ function
nor rules one out. The next mathematical problem is a
transform-level inequality that uses the exact modular kernel.

\paragraph{Provenance and DOI.}
Source repository:
\url{https://github.com/tonynewton79-web/RiemannHypothesis}.
Reserved, unpublished DOI:
\href{https://doi.org/10.5281/zenodo.23461586}{10.5281/zenodo.23461586}.
Neither peer review, independent novelty, nor a complete
Lean proof is presently asserted.

\section{A transform-level Laguerre attack}\label{sec:transformattack}
The Riemann hypothesis asks about zeros of the completed
xi-function, rather than the curvature of its kernel.
With the present Csordas--Planat--Sol\'e normalization,
write $K(u)=\Phi(u)$ and
\begin{equation}\label{eq:fourier-actual}
F(z)=\int_{\mathbb R}K(u)e^{izu}\,du
     =\frac14\,\xi\!\left(\frac12+\frac{iz}{2}\right).
\end{equation}
The factor $1/4$ and the half-frequency must be retained:
the first real zero of $F$ is near $z=28.269$, not $14.134$.
The scaling does not affect signs of generalized Laguerre
expressions.

For a sufficiently rapidly decreasing positive even kernel,
set
\begin{equation}
L_n[F](x)=\frac{(-1)^n}{(2n)!}
\sum_{j=0}^{2n}(-1)^j {2n\choose j}
F^{(j)}(x)F^{(2n-j)}(x).
\end{equation}
Differentiating the Fourier integral and symmetrizing
yields an exact correlation representation:
\begin{equation}\label{eq:all-order-corr}
L_n[F](x)=\frac{1}{(2n)!}
\iint_{\mathbb R^2}(u-v)^{2n}K(u)K(v)
\cos(x(u+v))\,du\,dv.
\end{equation}

\begin{proposition}[Central positivity and a local band]
Let $K$ be real, even, positive almost everywhere and have
the moments required in \eqref{eq:all-order-corr}.
Define the positive quantities
\begin{align}
A_n&=\frac{1}{(2n)!}\iint (u-v)^{2n}K(u)K(v)\,du\,dv,\\
B_n&=\frac{1}{(2n)!}\iint (u-v)^{2n}(u+v)^2K(u)K(v)\,du\,dv.
\end{align}
Then $L_n[F](0)=A_n>0$ and for every real $x$,
\begin{equation}\label{eq:local-positive}
L_n[F](x)\ge A_n-\frac{x^2}{2}B_n.
\end{equation}
In particular $L_n[F](x)>0$ for
$|x|<\sqrt{2A_n/B_n}$ whenever $B_n>0$.
\end{proposition}
\begin{proof}
Equation~\eqref{eq:all-order-corr} is the coefficient of
$y^{2n}$ in
$F(x+iy)F(x-iy)$, using symmetry under $u\leftrightarrow v$.
The strict inequality at $x=0$ follows from positivity of
the integrand away from the measure-zero diagonal.
Finally $\cos\theta\ge 1-\theta^2/2$ for every real $\theta$,
and integration against the nonnegative correlation weight
gives~\eqref{eq:local-positive}.
\end{proof}
This observation is a diagnostic obstruction, not an RH
proof: central positivity holds even for the explicit
counterexamples of the preceding section. The all-order
correlation construction has prior art in
Dimitrov--Xu~\cite{DX}.

For example, with $M_j=\int_{\mathbb R}u^j K(u)\,du$
and $M_{2j+1}=0$, direct expansion at order one gives
\begin{equation}\label{eq:n1moments}
A_1=M_0M_2,\qquad B_1=M_0M_4-M_2^2.
\end{equation}
The latter is nonnegative by Cauchy--Schwarz.

\section{Numerical exploration of the genuine transform}
The reproducible experiment
\texttt{certificates/generalized\_laguerre\_scan.py}
uses 65-decimal-digit \texttt{mpmath} quadrature for the
Jacobi-theta representation. It evaluates
$L_n[F](x)$ for $n=1,\ldots,5$ at
$x\in\{0,5,10,15,20,25,28,30\}$.
All 40 values were \emph{numerically} positive.
At the endpoints, $L_1[F](0)\simeq
1.78434856788\times 10^{-4}$ and
$L_1[F](30)\simeq 1.09904761087\times10^{-8}$.
Independent differentiation of the completed zeta function
in \eqref{eq:fourier-actual} at $x=0,10,28$ agrees to roughly
$10^{-60}$ relative precision for derivatives through order four.

The numerically estimated local positivity radii
$\sqrt{2A_n/B_n}$ for $n=1,\ldots,8$ are, respectively,
\[
9.8314,\ 10.3132,\ 10.7562,\ 11.1665,\
11.5491,\ 11.9081,\ 12.2466,\ 12.5674.
\]
These are exploratory evaluations, \emph{not} rigorously
enclosed moment integrals or proofs of positivity outside
the sampled points. Their growth suggests looking for
an $n$-uniform, noncentral correlation estimate, but no
such estimate is claimed here.

\section{Lean expansion and the remaining analytic bridge}
The new module \texttt{XiKernel/TransformMoment.lean}
proves a finite-sum central-moment positivity theorem for
every natural number $n$, the elementary quadratic gate,
and exact order-one moment identities. Its hypotheses
involve abstract strictly positive even moments;
the analytic identification with the real theta-kernel
integrals is \emph{not} formalized.

The module \texttt{XiKernel/RationalJet.lean}
defines a finite Taylor-jet recurrence and computes the
fifth iterate of rational 12-decimal approximants to the
eleven genuine theta Taylor coefficients. GitHub Lean CI
passed, but its \texttt{native\_decide} proof depends on
\texttt{Lean.ofReduceBool} and \texttt{Lean.trustCompiler},
as confirmed by an explicit \texttt{\#print axioms} audit.
It is compiler-assisted reflection, \emph{not} a small-kernel-only
certificate. This establishes only a rational sign for the
\emph{approximate} jet; it cannot be substituted for
Theorem~\ref{thm:five} without a formally checked
analytic coefficient enclosure and stability theorem.

To certify Theorem~\ref{thm:five} \emph{end to end in Lean},
three missing links must all be discharged:
(i) the theta-series differentiation and Taylor-coefficient
identities within Mathlib, including modular evenness;
(ii) formal rational enclosures for $\pi$, exponential
truncation, and every omitted $m\ge13$ summand;
(iii) sound interval propagation from these enclosures
to the actual fifth iterate and the strict negative bound.
The existing standard-library Python calculation supplies
a reproducible exact computational certificate but is not
a kernel-checked Lean derivation of (i)--(iii).
No \texttt{sorry}, fabricated axiom, or numerical
observation is to be represented as closing these steps.

\section{A conditional all-order noncentral band}
The numerical growth of the local radii suggests a more
structural criterion, independent of individual moment
calculations. It is presented as a conditional theorem,
\emph{not} as an established property of the actual theta
kernel.

\begin{proposition}[Radially increasing potential curvature]\label{prop:radial}
Let $K(u)=e^{-V(u)}$ be positive and even, with sufficient
decay to have all moments. Suppose $V$ is twice continuously
differentiable and $V''(u)$ is nondecreasing for $u\ge0$.
Then the ratios $B_n/A_n$ of the preceding section
are nonincreasing as $n$ increases:
\[
\frac{B_{n+1}}{A_{n+1}}\le
\frac{B_n}{A_n},\qquad n\ge0.
\]
Consequently the local radii $\sqrt{2A_n/B_n}$
are nondecreasing, and the single order-one bound,
if valid, gives $L_n[F](x)>0$ for every $n\ge1$
throughout $|x|<\sqrt{2A_1/B_1}$.
\end{proposition}
\begin{proof}
Set $d=u-v$, $w=u+v$. Up to a constant, the
joint density on the positive quadrant $d,w\ge0$ is
\[
g(d,w)=\exp\{-V((w+d)/2)-V((w-d)/2)\}.
\]
By evenness, $V''$ is even and therefore
\[
\partial_d\partial_w\log g
=\frac14\left[V''\!\left(\frac{w-d}{2}\right)
                  -V''\!\left(\frac{w+d}{2}\right)\right]
\le0,
\]
because $|w-d|\le w+d$. Thus for $d_2>d_1$
the likelihood ratio $g(d_2,w)/g(d_1,w)$
is nonincreasing in $w\ge0$. By the elementary
reverse-Chebyshev covariance inequality, the
conditional expectation $h(d)=
\mathbb E[w^2\mid |u-v|=d]$ is nonincreasing in $d$.

Under the tilted joint density proportional to
$|d|^{2n}g(|d|,|w|)$, one has
$B_n/A_n=\mathbb E_n[w^2]=\mathbb E_n[h(d)]$.
Passing from $n$ to $n+1$ multiplies the marginal
density of $d$ by the increasing weight $d^2$,
so this expectation of the decreasing function
$h$ cannot increase. This proves the ratio
inequality; the noncentral band follows from
\eqref{eq:local-positive}.
\end{proof}

\paragraph{Unproved theta-kernel gate.}
For the genuine Jacobi theta kernel, the required
global hypothesis is
\[
V'''(u)=-\big(\log\Phi(u)\big)'''\ge0
\qquad(u>0).
\]
High-precision numerical tests at
$u=0.0001,0.001,0.01,0.025,0.05,0.1,0.2,
0.4,0.8,1.5,2$ returned strictly positive values;
this is not a rigorous interval or an all-$u$
certificate. If proved, it would yield a
nontrivial all-order Laguerre positivity \emph{strip}
for the actual transform, but still not RH.
Checking prior art and finding a falsifying region
for the curvature condition are required before
making a claim of novelty.

\section{Endpoint curvature of the actual theta kernel}
\label{sec:curvatureendpoints}
The condition $V'''(u)>0$ for $u>0$,
with $V=-\log\Phi$, remains globally unresolved.
We isolate both endpoint regimes by exact analytic arguments.

\begin{proposition}[Modular-point cancellation]
Let
\[
\Phi_1(u)=\pi(2\pi e^{4u}-3)e^{5u-\pi e^{4u}},
\qquad V_1=-\log\Phi_1.
\]
For $y=\pi e^{4u}$,
\begin{align}
V_1'''(u)
&=\frac{64y}{(2y-3)^3}P(y),\\
P(y)&=(2y-3)^3-6(2y+3).
\end{align}
In particular $V_1'''(0)<0$, whereas
$V_1'''(u)>0$ whenever $y\ge7/2$.
The polynomial has exactly one real root
on $[3,\infty)$, lying between $13/4$
and $7/2$.
\end{proposition}
\begin{proof}
Writing $\log\Phi_1(u)=
\log\pi+\log(2y-3)+5u-y$
and applying $(4y\,d/dy)^3$ yields
the displayed expression. The polynomial is
\[
P(y)=8y^3-36y^2+42y-45.
\]
Its derivative $P'(y)=24(y-3/2)^2-12$
is positive for $y\ge3$. The bounds
$3<\pi<22/7<13/4$ and $P(13/4)=-113/8<0$
give $P(\pi)<0$, while $P(7/2)=4>0$.
For $y=t+7/2$ with $t\ge0$,
\[
P(y)=8t^3+48t^2+84t+4>0.
\]
Monotonicity and continuity give uniqueness.
\end{proof}

\begin{proposition}[Full-kernel endpoint positivity]
Write $s(t)=\Phi(\sqrt t)=a_0+a_1t+a_2t^2+O(t^3)$
with $a_0>0$. If $\mathcal T s(0)
=a_1^2-2a_0a_2>0$, then
\begin{equation}
V'''(u)=\frac{12\,\mathcal T s(0)}{s(0)^2}
\,u+O(u^3)\quad (u\to0^+).
\end{equation}
Hence $V'''(u)>0$ for all sufficiently
small $u>0$ for the genuine theta kernel.
Moreover $V'''(u)>0$ for all sufficiently
large $u$, without assuming a global
curvature hypothesis.
\end{proposition}
\begin{proof}
The theta functional equation makes $\Phi$
even near the modular fixed point.
Expanding
\[
-\log s(u^2)=-\log a_0-\frac{a_1}{a_0}u^2
-\left(\frac{a_2}{a_0}
-\frac{a_1^2}{2a_0^2}\right)u^4+O(u^6)
\]
and differentiating three times yields
the first assertion. Strict positivity of
$\mathcal T s(0)$ is available from the
earlier second-level kernel theorem
and independently from the exact interval
certificate in this work.

For large $u$ put $y=\pi e^{4u}$. The ratio
of the full sum to its first summand is
$1+R(y)$, where
\[
R(y)=\sum_{m\ge2}
m^2\frac{2m^2y-3}{2y-3}e^{-(m^2-1)y}.
\]
Termwise derivatives through order three
are uniformly dominated by a polynomial
in $y$ times $e^{-3y}$ for $y\ge y_0>3$;
this follows by a convergent geometric
majorant after bounding all factors by
polynomials in $m$ and $y$.
It follows that
$(4y\partial_y)^3\log(1+R(y))
=O(y^N e^{-3y})$ for some fixed $N$.
Since the first-summand curvature above
satisfies $V_1'''(u)=64y+O(1)$, we obtain
\[
V'''(u)=64\pi e^{4u}+O(1)
\qquad(u\to\infty),
\]
hence eventual strict positivity.
\end{proof}

The global problem is now concentrated on a
\emph{compact intermediate interval}. The
modular neighborhood is controlled through
the positive second-level defect; the distant
tail is controlled by first-summand dominance.
An explicit rational subdivision certificate
for the intermediate region, followed by an
analytic/Lean transfer, is still missing.
This endpoint theorem does not imply the
complete Laguerre hierarchy on the real line,
nor the Riemann hypothesis.

\section{Weighted-defect form of the global curvature question}
\label{sec:weightedcurvature}
There is a particularly direct equivalent form of the
remaining monotonicity condition, using the same variable
$t$ as the earlier Csordas--Dimitrov differential operator.

\begin{proposition}[Exact weighted-defect equivalence]
Let $s:(0,\infty)\to(0,\infty)$ be three times
continuously differentiable and put
\[
\Phi(u)=s(u^2),\quad V(u)=-\log\Phi(u),\quad
q(t)=\frac{\mathcal Ts(t)}{s(t)^2}
=-(\log s)''(t).
\]
For every $u>0$,
\begin{align}
V'''(u)&=12u\,q(u^2)+8u^3q'(u^2)\nonumber\\
&=8\left.\frac{d}{dt}
[t^{3/2}q(t)]\right|_{t=u^2}.
\label{eq:weightedeq}
\end{align}
Consequently
\[
V'''(u)\ge0\ \ (u>0)
\quad\Longleftrightarrow\quad
t^{3/2}\frac{\mathcal Ts(t)}{s(t)^2}
\text{ is nondecreasing on }t>0.
\]
\end{proposition}
\begin{proof}
Set $L(t)=\log s(t)$. Chain differentiation
gives
\[
V''(u)=-2L'(u^2)-4u^2L''(u^2)
\]
and hence
\[
V'''(u)=-12uL''(u^2)-8u^3L'''(u^2).
\]
Substituting $q=-L''$ proves the first identity,
and differentiating $t^{3/2}q(t)$ proves
the second. Since the change $u\mapsto u^2$
preserves the positive half-line, the final
equivalence follows.
\end{proof}

The known second-level theorem controls the
\emph{sign} of $\mathcal Ts$ but does not by
itself control the derivative of the weighted
quotient. Thus the new proof target is
\[
3q(t)+2tq'(t)\ge0,\qquad t>0.
\]
It is a one-dimensional inequality, involving
no Fourier integration, and could be tackled
through Jacobi's differential equation,
a compact interval certificate, or a direct
comparison inequality for $\mathcal Ts$ and
$s$. Even full success establishes only a
common finite Laguerre positivity band, not RH.

\section{Global radial-curvature certificate for the Xi kernel}
\label{sec:globalcurvature}

The monotonicity of the potential curvature for the
classical Riemann theta kernel is not new: Newman proved
in 1991 that $V'$ is convex on the positive half-line,
using a GHS-type argument from statistical mechanics
\cite{NewmanGHS}. Our aim here is different and narrower:
to give a reproducible quantitative verification by
outward-rounded rational intervals and an explicit
infinite-range estimate, while isolating the exact
bounds subsequently used in the Laguerre comparison.
In particular, this is not a claim of first discovery
of the underlying monotonic-curvature property.
The numerical and analytic certificate still requires
independent specialist review.

\begin{theorem}[Global monotone theta curvature]
\label{thm:globalcurvature}
For the classical Jacobi-theta kernel
\begin{align*}
\Phi(u)&=\sum_{m\ge1}\pi m^2
  \bigl(2\pi m^2e^{4u}-3\bigr)\\
&\qquad\cdot e^{5u-\pi m^2e^{4u}}.
\end{align*}
let $V(u)=-\log\Phi(u)$. Then
\begin{equation}\label{eq:globalthird}
\boxed{V'''(u)>0\quad\text{for every }u>0.}
\end{equation}
Consequently $V''$ is strictly increasing in
$|u|$. In the variable $t=u^2$, the function
\[
t^{3/2}\frac{\mathcal Ts(t)}{s(t)^2}
\quad (s(t)=\Phi(\sqrt t))
\]
is strictly increasing on $(0,\infty)$.
\end{theorem}

\paragraph{Exact reduction.}
Put $x=\pi e^{4u}$ and
\begin{align}\label{eq:Aj}
A_j(x)&=\sum_{m\ge1}(-1)^j m^{2+2j}
 \bigl(2m^2x-3-2j\bigr)\nonumber\\
&\qquad\cdot e^{-(m^2-1)x},
\quad 0\le j\le4.
\end{align}
Then
$\Phi(u)=\pi e^{5u-x}A_0(x)$ and
$A_j'(x)=A_j(x)+A_{j+1}(x)$ for $j\le3$.
All $A_j$ are absolutely convergent and
smooth for $x\ge\pi$. Since $A_0(x)>0$,
define the cubic numerator
\begin{align}
H(x)={}&-x^3
\bigl(A_3A_0^2-3A_1A_2A_0+2A_1^3\bigr)
\nonumber\\
&-3x^2\bigl(A_2A_0^2-A_1^2A_0\bigr)
-xA_1A_0^2.
\label{eq:H}
\end{align}
Three applications of $4x\partial_x$ to
$\log A_0(x)-x+5u+\log\pi$ give the
\emph{exact identity}
\begin{equation}\label{eq:thirdH}
V'''(u)=\frac{64H(x)}{A_0(x)^3}.
\end{equation}
The Jacobi functional equation makes $\Phi$
even about $u=0$, so $V'''(0)=0$ and
$H(\pi)=0$.

\paragraph{Finite interval portion.}
The standalone standard-library script
\texttt{certificates/xi\_global\_curvature\_exact.py}
works with common denominator $10^{65}$,
rounding every multiplication, reciprocal,
and integer division outward. A rigorous
Machin-series enclosure bounds $\pi$.
For $0\le z<1/4$, an 80-term Taylor
expansion encloses $e^z$, since
\[
0<e^z-\sum_{k=0}^{80}\frac{z^k}{k!}
 \le \frac{2}{4^{81}81!}<10^{-65}.
\]
Exponentials $e^{-y}$ are then enclosed by
\[
e^{-y}=\bigl(e^{y/4096}\bigr)^{-4096}.
\]
For every $x\in[\pi,6]$, truncating
\eqref{eq:Aj} at $m=8$ causes an error,
for each $0\le j\le4$, smaller than
\[
2\cdot23\cdot9^{12}\cdot2^{-240}
<10^{-58}.
\]
Indeed the $m$th absolute contribution
for $m\ge9$ is bounded by
$23m^{12}2^{-3(m^2-1)}$, and its
successive ratio is less than $1/2$.
The script explicitly enlarges all five
intervals by $10^{-58}$.

Using $A_j'=A_j+A_{j+1}$ and
forward automatic differentiation of
\eqref{eq:H}, $100$ rational interval cells
covering $[\pi,7/2]$ establish
\begin{equation}
\boxed{H'(x)>46\quad(\pi\le x\le7/2).}
\label{eq:Hprimebound}
\end{equation}
With $H(\pi)=0$, this yields $H(x)>0$
for $\pi<x\le7/2$. Independently,
$400$ rational interval cells on $[7/2,6]$
establish
\begin{equation}
\boxed{H(x)>26\quad(7/2\le x\le6).}
\label{eq:Hcompactbound}
\end{equation}
These are outward enclosures on whole
intervals, not tests only at mesh nodes.

\paragraph{Analytic infinite tail.}
For $x\ge6$, let
\[
a_j=(-1)^j(2x-3-2j),\qquad E_j=A_j-a_j.
\]
For $0\le j\le3$ and $m\ge2$,
the $m$th summand of $E_j$ has magnitude
at most $3xm^{10}e^{-(m^2-1)x}$.
Successive terms have ratio bounded by
$(3/2)^{10}e^{-30}<1/2$.
Consequently
\begin{equation}\label{eq:Etail}
|E_j(x)|\le6144x e^{-3x}=:\epsilon(x)x,
\qquad 0\le j\le3.
\end{equation}
The first-summand numerator, a purely
algebraic simplification, is
\begin{equation}\label{eq:H0}
H_0(x)=x(8x^3-36x^2+42x-45)
>\frac74x^4
\end{equation}
for $x\ge6$. The latter strict inequality
follows, for $t=x-6\ge0$, from
\begin{align*}
H_0(x)-\frac74x^4
={}&\frac{25}{4}t^4+114t^3+744t^2\\
&+1971t+1566>0.
\end{align*}
Along every straight-line interpolation
between $(a_0,\ldots,a_3)$ and
$(A_0,\ldots,A_3)$, we have
$|A_j|\le3x$ because $\epsilon(x)<1$.
Bounding the four first partial derivatives
of the cubic \eqref{eq:H} by the
multivariate mean-value inequality yields
\begin{equation}\label{eq:Herror}
|H-H_0|\le \epsilon(x)x^4
(162x^2+162x+27).
\end{equation}
The function
$e^{-3x}(162x^2+162x+27)$ is
strictly decreasing for $x\ge6$.
Finally the elementary rational Taylor bound
$e^3>20$ gives $e^{18}>20^6=64{,}000{,}000$;
hence for every $x\ge6$,
\[
\frac{|H-H_0|}{x^4}
<\frac{6144\cdot6831}{64{,}000{,}000}
<\frac74.
\]
Equations \eqref{eq:H0}--\eqref{eq:Herror}
give $H(x)>0$ on $[6,\infty)$.
Combining the three ranges with
\eqref{eq:thirdH} proves
Theorem~\ref{thm:globalcurvature}.

\begin{proposition}[Unconditional common Laguerre band]
\label{prop:commonlaguerre}
Let $F(z)=\int_\mathbb{R}\Phi(u)e^{izu}\,du$,
and let $A_n,B_n$ be the positive
correlation moments of
Section~\ref{sec:transformattack}.
For every integer $n\ge1$,
\[
\frac{B_n}{A_n}\le\frac{B_1}{A_1}.
\]
In particular,
\[
\boxed{L_n[F](x)>0\quad
\text{for all }n\ge1,\quad
|x|<\sqrt{2A_1/B_1}.}
\]
\end{proposition}
\begin{proof}
Theorem~\ref{thm:globalcurvature}
establishes the radial monotonicity
hypothesis of
Proposition~\ref{prop:radial}.
The reverse-likelihood-ratio argument
therefore proves the moment-ratio
inequality. Applying
$L_n[F](x)\ge A_n-\tfrac12x^2B_n$
finishes the proof.
\end{proof}

\paragraph{Certification and limitations.}
The two finite inequalities above are
established by exact integer operations and
explicit analytic theta/exponential tails;
the entire proof is independently auditable
from the standalone Python 3 script.
The algebraic identities and the rational
infinite-tail margin are included in the
Lean module
\texttt{XiKernel/GlobalCurvatureAlgebra.lean}.
The \emph{full} theta-series positivity
statement and the interval verifier have
not yet been formally imported and proved
within Lean's kernel.
Neither the global radial-curvature theorem
nor its common finite Fourier-positivity
band implies the Riemann hypothesis.
Newman's prior theorem \cite{NewmanGHS} establishes
the principal qualitative convexity assertion.
Any possible contribution of this section is the
alternative constructive certificate, its strict
positive rational margins on the stated domains,
and the way those estimates are applied here;
their novelty has not been independently established.
Independent specialist review remains outstanding.

\section{An exact falsification of the curvature-to-real-zero inference}
\label{sec:curvatureobstruction}

The global radial-curvature theorem for the genuine Jacobi-theta
kernel provides a common \emph{bounded} Fourier-side Laguerre
positivity interval. It cannot on its own be extended to
the entire real Fourier axis. The following exact example
makes this logical boundary decisive.

\begin{theorem}[Gaussian--hyperbolic counterexample]
\label{thm:coshcounter}
Let
\[
K(u)=e^{-u^2/2}\left(1+\frac12\cosh u\right),
\qquad V(u)=-\log K(u).
\]
Then $K$ is positive, even, entire, and rapidly decreasing
on the real axis, with
\begin{align}
V''(u)
&=\frac{\cosh^2u+2\cosh u+3}
        {(\cosh u+2)^2}\ge\frac23,\\
V'''(u)
&=\frac{2(\cosh u-1)\sinh u}
        {(\cosh u+2)^3}>0\quad(u>0).
\label{eq:coshcurv}
\end{align}
Nevertheless its Fourier transform is
\begin{equation}
F(z)=\int_{\mathbb R}K(u)e^{izu}\,du
=\sqrt{2\pi}e^{-z^2/2}(1+d\cos z),
\qquad d=\frac{e^{1/2}}2,
\label{eq:coshfourier}
\end{equation}
where $1/2<d<1$. Thus $F(x)>0$ for \emph{every}
real $x$, but $F$ has nonreal zeros
\begin{equation}
z=(2k+1)\pi\pm i\operatorname{arcosh}(1/d),
\qquad k\in\mathbb Z.
\end{equation}
Its first Fourier-side Laguerre inequality is
already violated at the explicit real point $x=\pi$:
\begin{equation}\label{eq:coshlaguerre}
\boxed{
L_1[F](\pi)
=2\pi e^{-\pi^2}(1-d)(1-2d)<0.
}
\end{equation}
\end{theorem}

\begin{proof}
Put $C=\cosh u$, $S=\sinh u$. The elementary identities
$C^2-S^2=1$, $C'=S$, and $S'=C$ yield
\[
V''=1-\frac{2C+1}{(C+2)^2}
\quad\text{and}\quad
V'''=\frac{2(C-1)S}{(C+2)^3}.
\]
The strong convexity bound is exactly
\[
3(C^2+2C+3)-2(C+2)^2
=(C-1)^2\ge0.
\]
For $u>0$, $C>1$ and $S>0$, proving
strict radial curvature.

The standard Gaussian Fourier integral
$\int_{\mathbb R}e^{-u^2/2}e^{izu}du
=\sqrt{2\pi}e^{-z^2/2}$ holds for complex $z$
by analytic continuation.
The hyperbolic factor is
$\cosh u=(e^u+e^{-u})/2$; shifting the Fourier
argument by $\mp i$ and completing squares
yields \eqref{eq:coshfourier}.
Elementary exponential series show
$1<e^{1/2}<2$, hence $1/2<d<1$.
For real $x$, $1+d\cos x\ge1-d>0$.
On the other hand, at
$z=(2k+1)\pi+iy$ one has
$\cos z=-\cosh y$, and
$\cosh y=1/d>1$ has a unique $y>0$,
proving the explicit nonreal zero pair.

Finally differentiating
$F(z)=C_0e^{-z^2/2}H(z)$,
$H(z)=1+d\cos z$, gives
\[
(F')^2-FF''
=C_0^2e^{-z^2}
\bigl[H^2+(H')^2-HH''\bigr].
\]
Substituting $z=\pi$, where
$H=1-d$, $H'=0$, $H''=d$, gives
\eqref{eq:coshlaguerre}. Every sign follows
from $1/2<d<1$.
\end{proof}

\begin{proposition}[A sharp zero-location transition
inside the radial-curvature class]\label{prop:coshfamily}
For any $c\ge1/2$ define
\[
K_c(u)=e^{-u^2/2}(1+c\cosh u),\qquad
d_c=c\,e^{1/2}.
\]
All these kernels have $(-\log K_c)'''(u)>0$
for $u>0$, and their potentials are
strictly convex. Their Fourier transforms are
\[
F_c(z)=\sqrt{2\pi}e^{-z^2/2}(1+d_c\cos z).
\]
The transform has only real zeros exactly when
$d_c\ge1$, and nonreal zeros when $0<d_c<1$.
Thus the continuous curvature-positive family
passes through a real-zero transition at
$c=e^{-1/2}$, a parameter strictly greater than $1/2$.
\end{proposition}
\begin{proof}
Write $M=1+c\cosh u$ and $C=\cosh u$.
Direct differentiation yields
\[
(\log M)'''=
\frac{c\sinh u(1-cC-2c^2)}
{(1+cC)^3}.
\]
For $c\ge1/2$ and $u>0$, one has
$cC+2c^2>1$, so $(-\log K_c)'''>0$.
Moreover
\[
(\log M)''=
\frac{cC+c^2}{(1+cC)^2}<1,
\]
since the difference between denominator
and numerator is $1+cC+c^2(C^2-1)>0$.
Thus $(-\log K_c)''>0$.
Gaussian integration gives the transform.
If $d_c<1$, $\cos z=-1/d_c<-1$ has
nonreal solutions and no real solutions.
If $d_c\ge1$, solving
$\cos z=-1/d_c\in[-1,0)$ shows all
solutions are real (the boundary case
$d_c=1$ has double real zeros).
\end{proof}

\paragraph{Implications for RH.}
The counterexample proves that positivity of
$V'''$, even coupled to uniform strong log-concavity,
cannot imply either $L_1[F](x)\ge0$ for all real $x$
or the real-zero property. This does not weaken
Theorem~\ref{thm:globalcurvature} for the \emph{specific}
theta kernel; it establishes that a successful RH proof
must exploit additional theta/modular information
\emph{at the Fourier-transform level}.
A precise future target is a signed-cosine
correlation inequality
\[
\int_{\mathbb R}
G_n(w)\cos(xw)\,dw\ge0\qquad
(\forall n\ge1,\ \forall x\in\mathbb R),
\]
where
$G_n(w)=\int(2u-w)^{2n}\Phi(u)\Phi(w-u)\,du$
with the positive factorial normalization
specified earlier. The mere fact $G_n(w)\ge0$
does not provide this Fourier-positivity statement.
This distinction is consistent with the earlier
Wronskian/correlation framework of
Dimitrov--Xu~\cite{DX}.

\section{An explicit arbitrary-depth obstruction to finite Laguerre testing}
\label{sec:finitedepth}
The preceding Gaussian--hyperbolic example
shows that radial curvature does not force even
the first Fourier-side Laguerre inequality.
There is a stronger structural obstruction:
\emph{no fixed finite number of generalized
Laguerre inequalities, even when valid for
every real Fourier variable, implies the
real-zero property.}

\begin{theorem}[Finite-depth obstruction]
\label{thm:finitedepth}
For every integer $N\ge1$ put
\[
d_N=2^{-(N+2)},\qquad
K_N(u)=e^{-u^2/2}
\bigl(1+d_N e^{-1/2}\cosh u\bigr)
\]
and $F_N(z)=\int_{\mathbb R}K_N(u)e^{izu}\,du$.
Then $K_N$ is positive, even, rapidly
decreasing and
\begin{equation}
F_N(z)=\sqrt{2\pi}e^{-z^2/2}
\bigl(1+d_N\cos z\bigr).
\label{eq:finiteF}
\end{equation}
For every $1\le n\le N$ and $x\in\mathbb R$,
\begin{equation}
\boxed{
L_n[F_N](x)\ge
\frac{\pi e^{-x^2}}{n!}>0.
}
\label{eq:finitelb}
\end{equation}
Nevertheless $F_N$ has nonreal zeros
\begin{equation}
z=(2k+1)\pi\pm
i\operatorname{arcosh}(1/d_N),
\qquad k\in\mathbb Z.
\label{eq:finitezero}
\end{equation}
\end{theorem}

\begin{proof}
Gaussian Fourier integration gives
\eqref{eq:finiteF} and
\eqref{eq:finitezero} follows as in
Theorem~\ref{thm:coshcounter}.
The generalized Laguerre coefficients
are determined by the generating function
\begin{align*}
&F_N(x+iy)F_N(x-iy)\\
&\quad=2\pi e^{-x^2}e^{y^2}
\bigl[1+2d_N\cos x\cosh y\\
&\qquad\quad+\tfrac12d_N^2(\cosh 2y+\cos 2x)\bigr].
\end{align*}
Writing $L_n[F_N](x)$ for the
coefficient of $y^{2n}$, set
\begin{align*}
a_n&=\frac1{n!},\\
b_n&=\sum_{j=0}^n
  \frac1{(n-j)!(2j)!},\\
c_n&=\sum_{j=0}^n
  \frac{4^j}{(n-j)!(2j)!}.
\end{align*}
Then
\begin{align}\label{eq:finitecoefs}
\frac{L_n[F_N](x)}{2\pi e^{-x^2}}
&=a_n+2d_Nb_n\cos x\nonumber\\
&\quad+\frac{d_N^2}{2}
\left(c_n+a_n\cos 2x\right).
\end{align}
The last parenthesis is nonnegative
because $c_n\ge a_n$. Next
\[
\frac{b_n}{a_n}
=\sum_{j=0}^n
\binom nj\frac{j!}{(2j)!}
\le\sum_{j=0}^n\binom nj=2^n.
\]
Hence for $1\le n\le N$,
\begin{align*}
\frac{L_n[F_N](x)}{2\pi e^{-x^2}}
&\ge a_n-2d_Nb_n\\
&\ge a_n(1-2d_N2^n)\\
&\ge a_n/2>0,
\end{align*}
which proves \eqref{eq:finitelb}.
All bounds hold uniformly for
$x\in\mathbb R$, not just on compact sets.
\end{proof}

\begin{proposition}[First-order positivity hides
a second-order failure]
\label{prop:firstsecond}
With $d=2/5$, set
\[
F(z)=\sqrt{2\pi}\,e^{-z^2/2}
\left(1+\frac25\cos z\right).
\]
Then $F(x)>0$ and $L_1[F](x)>0$
for every real $x$, but
\begin{equation}
\boxed{L_2[F](\pi)=
-\frac{2\pi e^{-\pi^2}}{25}<0.}
\end{equation}
\end{proposition}
\begin{proof}
With $t=\cos x\in[-1,1]$,
\begin{align*}
\frac{L_1[F](x)}{2\pi e^{-x^2}}
&=1+d^2+3dt+d^2t^2\\
&\ge (1-d)(1-2d)=\frac3{25}>0.
\end{align*}
Indeed, for $d=2/5$ the difference
from $3/25$ is
$(t+1)(4t+26)/25\ge0$.
For $L_2$ take $x=\pi$.
Since $\cos(\pi\pm iy)=-\cosh y$,
\[
F(\pi+iy)F(\pi-iy)
=2\pi e^{-\pi^2}e^{y^2}
(1-d\cosh y)^2.
\]
Its $y^4$ coefficient is
\[
\frac12-\frac{25}{12}d+
\frac{11}{6}d^2
=\frac{22d^2-25d+6}{12}.
\]
At $d=2/5$ this is $-1/25$,
as required.
\end{proof}

\paragraph{Exact tests and formalization scope.}
The standard-library script
\texttt{certificates/finite\_depth\_laguerre\_exact.py}
checks the coefficient formulas by rational
arithmetic through $N=100$, as well as
the exact $d=2/5$ separation.
The Lean module
\texttt{XiKernel/FiniteDepthBarrier.lean}
formally proves the associated
real-algebraic positivity gate and low-order
coefficient identities. A complete Lean
formalization of the Fourier--Gaussian
integral and all-$N$ binomial calculation
is not yet claimed.

\paragraph{Consequence for the Riemann hypothesis.}
This theorem proves, within a positive
analytic kernel class, that even globally
positive inequalities
$L_1,\ldots,L_N$ for any fixed $N$
cannot prove all zeros real.
The point is logical and methodological,
not an assertion of novelty;
Laguerre--P\'olya characterizations are
classical and the correlation viewpoint
has established prior art
\cite{DX}.
For the genuine Riemann $\Xi$, the
research target must therefore be a
uniform argument in \emph{all} orders
$n$, or a materially stronger zero
exclusion principle that explicitly uses
Jacobi/theta arithmetic structure.

\section{A stronger obstruction: nonreal zeros arbitrarily near the real axis}
\label{sec:nearreal}
Theorem~\ref{thm:finitedepth} exhibits
nonreal zeros that recede into the complex plane
as $N$ grows. That distance restriction is
not essential. The Gaussian parameter can be
chosen to conceal \emph{arbitrarily near-real}
complex zeros from every prescribed finite
number of global Laguerre tests.

\begin{theorem}[Near-real finite-depth barrier]
\label{thm:nearreal}
For every integer $N\ge1$ and every
$\varepsilon>0$ there exists a positive,
even, rapidly decreasing real kernel $K$
whose Fourier transform $F$ has nonreal
zeros at distance less than $\varepsilon$
from the real axis, although
\[
L_n[F](x)>0
\qquad\text{for every real }x,\quad
1\le n\le N.
\]
More explicitly, choose any $d\in(0,1)$,
set
\begin{equation}
A=\frac{16N}{(1-d)^2},\qquad
K(u)=e^{-u^2/(2A)}
\left[1+d\,e^{-1/(2A)}\cosh(u/A)\right],
\label{eq:nearK}
\end{equation}
and define $F(z)=\int K(u)e^{izu}\,du$.
Then
\[
F(z)=\sqrt{2\pi A}\,e^{-Az^2/2}
(1+d\cos z)
\]
and for all real $x$ and $1\le n\le N$
\begin{equation}
\boxed{
L_n[F](x)\ge
\frac{3\pi A}{2}
\,e^{-Ax^2}\,
\frac{(1-d)^2A^n}{n!}>0.
}
\label{eq:nearbound}
\end{equation}
The nonreal zeros are precisely
\[
(2k+1)\pi\pm i\operatorname{arcosh}(1/d).
\]
As $d\uparrow1$, their imaginary parts
tend to zero.
\end{theorem}

\begin{proof}
The inverse Fourier formula follows by
completing the square in the Gaussian
integral, with the $\cosh(u/A)$ term
corresponding to shifts $z\mapsto z\pm i/A$.
The resulting kernel is positive and
rapidly decaying for real $u$.

Fix $1\le n\le N$, and put
\begin{align*}
a&=\frac{A^n}{n!},\\
b&=\sum_{j=1}^n
\frac{A^{n-j}}{(n-j)!(2j)!},\\
c&=\sum_{j=1}^n
\frac{4^jA^{n-j}}{(n-j)!(2j)!}.
\end{align*}
The coefficient of $y^{2n}$ in
$e^{Ay^2}(1+d\cos(x+iy))
(1+d\cos(x-iy))$ is
\begin{equation}
a(1+d\cos x)^2
+2d\,b\cos x+\frac{d^2}{2}c.
\label{eq:nearcoefficient}
\end{equation}
Here $c\ge0$. The other two terms have
the lower bound
\[
a(1-d)^2-2db
\]
because $-1\le\cos x\le1$.
To estimate $b$, use
\[
\frac ba=
\sum_{j=1}^n
\frac{n!}{(n-j)!(2j)!A^j}
\le\sum_{j=1}^n\left(\frac nA\right)^j.
\]
Since $A\ge16N\ge2n$, the
geometric series is bounded by
\[
\frac ba\le
\frac{n/A}{1-n/A}\le\frac{2N}{A}.
\]
Therefore
\[
2db\le\frac{4dN}{A}a
=\frac{d(1-d)^2}{4}a.
\]
Using $d<1$ gives
\[
a(1-d)^2-2db
\ge\frac34a(1-d)^2>0.
\]
Multiplying by the positive Gaussian
factor $2\pi A\,e^{-Ax^2}$ proves
\eqref{eq:nearbound}.

Finally, take for instance
$d=1-m^{-2}$ with integer $m\ge2$.
Then $\cosh y=1/d$ at the nonreal
zero height $y>0$, and
$1/d=1+1/(m^2-1)$.
Since $\cosh\varepsilon\ge
1+\varepsilon^2/2$, any
$m$ satisfying
$m^2-1>2/\varepsilon^2$
gives $0<y<\varepsilon$.
This explicitly establishes the
arbitrary-nearness statement.
\end{proof}

\paragraph{Certification.}
The standard-library script
\texttt{certificates/near\_real\_finite\_depth\_exact.py}
checks the exact rational coefficient
bounds for several pairs $(N,m)$,
including a height threshold below
$10^{-3}$. The companion Lean module
\texttt{XiKernel/NearRealFiniteDepth.lean}
proves the underlying algebraic
positivity gate under the explicit
coefficient hypotheses; the
all-$N$ binomial comparison and
Gaussian integral are given above
as exact paper proofs rather than
end-to-end Lean theorems.

\paragraph{Research significance.}
Even global validity of any
\emph{finite initial block}
of generalized Laguerre inequalities
cannot detect nonreal zeros arbitrarily
near the real axis in this positive
Fourier-kernel class.
The issue is not the distance of a
putative off-line zero, but uniform
control of all Laguerre orders.
This does not decide any zero of
the genuine Riemann $\Xi$ function.

\section{Lean connection to the actual Riemann hypothesis}
\label{sec:leanxi}
The Lean formalization now reaches the \emph{genuine}
completed zeta function rather than just a symbolic
Gaussian model. Mathlib's
\texttt{completedRiemannZeta\(\_0\)}
represents the entire pole-subtracted completion
$\Lambda_0(s)$, satisfying
\[
\Lambda_0(1-s)=\Lambda_0(s),\qquad
\Lambda_0(s)=\Lambda(s)+\frac1s+\frac1{1-s}
\]
off the poles of $\Lambda(s)=\pi^{-s/2}\Gamma(s/2)\zeta(s)$.
Thus the identity
\begin{equation}
\xi(s)=\frac{1+s(s-1)\Lambda_0(s)}{2}
\label{eq:xi_mathi}
\end{equation}
defines an entire function with the correct extension
at $s=0$ and $s=1$. Its critical-line reparameterization is
\[
\widehat\Xi(z)=\xi\!\left(\frac12+\frac{iz}{2}\right).
\]
The Lean modules \texttt{XiKernel/EntireXi.lean} and
\texttt{XiKernel/AxiomAudit.lean} formally check
the $\xi$-functional equation, $\xi(0)=\xi(1)=1/2$,
complex differentiability, and the evenness and
entireness of $\widehat\Xi$.
The exact convention used earlier is
$F(z)=\widehat\Xi(z)/4$.

The proof target is to establish the
\emph{real-zero property}
\[
\widehat\Xi(z)=0\quad\Longrightarrow\quad
\operatorname{Im}z=0.
\]
To turn this into Mathlib's native
\texttt{RiemannHypothesis}, the formalization
must additionally relate zeros of $\xi$ to the
nontrivial zeros of $\zeta$ using the completed
zeta representation, with the trivial zeros
and poles treated explicitly. That zero
correspondence is a separate outstanding
proof obligation. The present Lean results
do not establish the real-zero property.

\section{Reproducibility and mathematical scope}
\label{sec:verification-release}
The numerical claims in this paper arise from three
different forms of calculation, and their logical roles
should be kept separate. The fifth-iterate calculation
for the genuine theta kernel uses outward rational
integer intervals, a Machin-series enclosure for
$\pi$, and explicit control of the omitted theta
summands. Its certificate is reproduced by
\path{certificates/xi_fifth_iterate_exact.py}.
The Gaussian-polynomial identities, the
large-frequency truncation coefficients and
the rational bounds on the exponential theta
tail can be checked directly from the
recurrences given in the proofs.

The radial monotone-curvature property of the classical
theta potential was previously proved by Newman in 1991
using a GHS inequality \cite{NewmanGHS}. The certificate
provided here is therefore an alternative quantitative
verification, not the discovery of that property.
The radial-curvature certificate is a
computer-assisted analytic result.
Its proof separates a compact interval,
certified by 500 outward-rounded rational
cells, from a distant region controlled
by an explicit infinite-tail majorant.
The implementation and the relevant
polynomial identities appear in
\path{certificates/xi_global_curvature_exact.py}.
Although the arithmetic checks are exact,
an independent mathematical review of
the complete tail comparison and the
placement of this theorem within earlier
theta-kernel literature remains appropriate.

The four finite zero-count theorems certify
distinct disks. The first three use a two-mode
Fourier approximation; the fourth requires
three modes and tighter local derivative bounds.
Their proofs reduce to enclosures of
three Taylor coefficients and one
positive remainder integral. Every
coefficient is computed at two
different composite-Simpson resolutions;
the stated fourth-derivative bounds
and the omitted $r>2$ integral tails
are handled analytically.
The integration nodes use
\texttt{mpmath.iv} version 1.3.0 for
interval enclosures of $\pi$,
exponentials and trigonometric values,
which are subsequently converted to
exact rational endpoints.
Thus the computational trust boundary
includes that library's transcendental
interval implementation as well as
the elementary analytic estimates.
The scripts, available in the
\texttt{certificates} directory,
may be run independently of the
LaTeX build. The corresponding tests
and their results are recorded
in the public repository's
\texttt{RESULTS.md} file.

The Lean 4.24.0/Mathlib development
certifies exact algebraic identities,
moment-positivity statements, rational
comparison gates, and the functional
symmetry of an entire form of
the classical completed zeta function.
For these small algebraic theorems,
the recorded axiom inspection reports
the ordinary Mathlib axioms
\texttt{propext}, \texttt{Classical.choice},
and \texttt{Quot.sound}.
The large rational-jet test instead
uses \texttt{native\_decide}, which
also relies on compiler-trust axioms;
it is not a small-kernel-only check.
Neither the interval-Simpson
calculations nor the full
Jacobi transformation, Fourier
integrals and Rouch\'e transfers
have been formalized end to end
in Lean.

The existence and simplicity of
the four zeros certified here
are classical results. Indeed,
Platt and Trudgian verified the
critical-line location and
simplicity of all nontrivial
zeta zeros up to height
$3\cdot10^{12}$ \cite{PlattTrudgian}.
The present certificates therefore
make no claim of a new zeta-zero
location or an improvement in
verified height. Their value is
to test the particular theta-Fourier
approximation and its explicit
Rouch\'e error bounds.

The open difficulty remains global.
The explicit theta truncation estimate
can be made arbitrarily small on the
complex strip, with an adaptive
number of modes growing as
$O(\sqrt{\log R})$ for prescribed
polynomial accuracy at height $R$.
What is unavailable is a
height-uniform strict lower bound
on the approximating transform
over contours that would exclude
off-critical zeros. Nor do the
finite Laguerre tests supply that
missing lower bound, as the
counterexamples in this paper show.
Consequently neither the four
bounded zero counts nor the
formal algebraic results prove
the Riemann hypothesis.

The complete source, reproducibility
instructions and verification record
are maintained at
\url{https://github.com/tonynewton79-web/RiemannHypothesis}.

\begin{thebibliography}{9}
\bibitem{CD} G. Csordas and D. Dimitrov,
Conjectures and theorems in the theory of entire functions,
\emph{Numer. Algorithms} 25 (2000), 109--122.
\bibitem{Cs} G. Csordas, Fourier transforms of positive definite kernels
and the Riemann $\xi$-function, \emph{CMFT} 15 (2015), 373--391;
\href{https://arxiv.org/abs/1309.0055}{arXiv:1309.0055}.
\bibitem{PS} M. Planat and P. Sol\'e,
\href{https://arxiv.org/abs/2608.19160}{arXiv:2608.19160} (2026).
\bibitem{LM} Y. Liu and J. Mao,
\href{https://arxiv.org/abs/2610.04972}{arXiv:2610.04972} (2026).
\bibitem{DX} D. Dimitrov and Y. Xu,
\href{https://arxiv.org/abs/1606.05011}{arXiv:1606.05011} (2016).
\end{thebibliography}
\end{document}
